\documentclass[11pt]{article}
\usepackage[margin=1in]{geometry}
\usepackage{amsmath,amssymb,amsthm,mathtools}
\usepackage{graphicx}
\usepackage{booktabs}
\usepackage{multirow}
\usepackage{xcolor}
\usepackage{hyperref}
\usepackage[numbers,sort&compress]{natbib}
\usepackage{microtype}
\usepackage{caption}
\usepackage{subcaption}
\usepackage{algorithm}
\usepackage{algpseudocode}
\usepackage{enumitem}

\hypersetup{
  colorlinks=true, linkcolor=blue!55!black, citecolor=blue!55!black,
  urlcolor=blue!55!black, pdftitle={Counterfactual Spatiotemporal Diffusion for Basketball}
}

\newcommand{\R}{\mathbb{R}}
\newcommand{\E}{\mathbb{E}}
\newcommand{\dd}{\mathrm{d}}
\newcommand{\loss}{\mathcal{L}}

\title{\bf Counterfactual Spatiotemporal Diffusion for Basketball:\\[2pt]
\large Stress-Testing Defensive Schemes via Multi-Agent Trajectory Generation}
\author{Anonymous submission}
\date{September 2026}

\begin{document}
\maketitle

\begin{abstract}
Defensive quality in basketball is ordinarily assessed retrospectively: a
coverage either succeeded or failed on the possessions where it was actually
deployed, confounding scheme quality with offensive execution, personnel, and
luck. We propose an alternative evaluation regime built on conditional
generative simulation. A score-based diffusion model learns the joint
distribution of half-court pick-and-roll (PnR) possessions, represented as
tensors $X \in \R^{T \times 11 \times 4}$ (ten players and the ball; positions
and velocities) in a rim-centric court frame, conditioned on the defensive
scheme, a handler-type flag, and the observed opening state. Because the model
is conditioned on intervention targets, sampling from it while \emph{holding
the opening state fixed} yields matched counterfactual futures: the same
offensive setup simulated under drop, switch, and blitz coverages, or under a
star-versus-role handler substitution. We contribute (i) a validated ingestion
and tensorization pipeline for the public 2015--16 NBA SportVU corpus,
including team inference, rim-frame normalization, tracker-noise rejection,
and geometric coverage labeling; (ii) a temporal U-Net diffusion architecture
with cross-attention conditioning and classifier-free guidance; (iii)
physics-informed training regularization and sampling-time gradient guidance
that we show are necessary for stable kinematically plausible generation; and
(iv) a vulnerability protocol over roller openness, paint pressure, and an
expected-possession-value proxy, together with a statistical-power analysis
showing that distribution-level realism metrics are session-stable at 16
samples per branch while branch-level tactical contrasts require samples in
the hundreds --- a negative result we regard as a necessary specification for
any use of generative counterfactuals in sports analytics. The complete
pipeline runs end-to-end on a single GPU and is released with a Colab
notebook, unit tests, and a behavioral simulator that doubles as an
interpretable baseline and counterfactual ground truth.
\end{abstract}

\section{Introduction}
\label{sec:intro}

Optical tracking transformed basketball analytics: every player and the ball
are localized at 25\,Hz, and a decade of research has converted those streams
into spatial metrics --- spacing \citep{skinner2015pitch}, court-occupation
intensities \citep{miller2014factorized}, possession-value processes
\citep{cervone2016multivariate,yamamoto2022epv}, and shot-probability
surfaces. Yet these tools share a structural limitation: they are
\emph{observational}. A drop coverage is evaluated only on the possessions
where a coach actually called drop, against the offenses that actually faced
it, with outcomes confounded by execution quality and matchup luck. Coaches
answer scheme questions with scrimmages and film; analysts answer them with
stratified observational statistics that grow unstable as the conditioning
stack deepens.

Generative simulation offers a different epistemic route. If a model captures
the conditional distribution $p(\text{futures} \mid \text{opening state},
\text{scheme}, \text{personnel})$, then a scheme comparison becomes a
\emph{controlled intervention}: fix one opening state, mutate the scheme
label, sample $M$ futures from each branch, and compare outcome
distributions. Confounds are broken by construction --- the offense never
``knows'' which coverage it will face, because the offense's opening state is
identical across branches. This is the counterfactual logic of randomized
experiments imported to spatial tactics, and it is only available if the
generative model is simultaneously (a) distributionally faithful, (b)
conditionable on tactical semantics, and (c) physically coherent.

We develop this program for the half-court pick-and-roll, the most frequent
and most scheme-sensitive action in modern basketball. Our contributions:

\begin{enumerate}[leftmargin=*,itemsep=2pt]
\item \textbf{A validated ingestion path for real tracking data.} We parse the
  raw 2015--16 SportVU archives (a format notably divergent from common
  descriptions of it), reconstruct contiguous possessions via epoch timestamps
  and entity signatures, infer teams and the attacked hoop, map coordinates
  into a rim-centric frame, reject tracker artifacts with median filtering and
  Savitzky--Golay smoothing, detect ball screens geometrically, and label
  coverages by defensive geometry (\S\ref{sec:data}). Every step is unit-tested
  and validated end-to-end on a downloaded public game.
\item \textbf{A conditional diffusion model for multi-agent possessions.} A
  temporal U-Net over the $T$ axis, with agent--feature grids folded into
  channels and cross-attention conditioning on scheme, handler type, and the
  anchor frame; cosine-schedule DDPM training, DDIM sampling, and
  classifier-free guidance through learned null tokens (\S\ref{sec:method}).
\item \textbf{Physics-informed training and sampling.} We show that naive
  physics regularization destabilizes diffusion training through the
  $1/dt^2$ amplification of finite-difference acceleration on noisy
  reconstructions, and derive a stable recipe: dimensionless relative-violation
  hinges evaluated on low-noise probes, $\bar\alpha_k$-weighted, with the bulk
  of constraint enforcement applied as gradient guidance on intermediate $X_0$
  predictions during sampling (\S\ref{sec:physics}).
\item \textbf{A vulnerability protocol and its power analysis.} Matched-branch
  counterfactuals (scheme mutation, handler swap) scored by roller openness,
  paint pressure, and an expected-possession-value proxy; a transparent
  behavioral simulator supplies interpretable baseline contrasts and label
  balance. Replicated GPU sessions establish that realism metrics are stable
  across runs while 16-sample branch contrasts are not statistically
  resolvable, and we derive the sampling budget needed to make them so
  (\S\ref{sec:experiments}).
\item \textbf{Reproducibility.} The repository executes the entire study ---
  archive download through evaluation report --- via one command on CPU and
  one Colab notebook on GPU, including 18 unit tests and animated court-space
  renders of generated possessions.
\end{enumerate}

The paper is organized as follows. \S\ref{sec:related} situates the work.
\S\ref{sec:problem} fixes notation. \S\ref{sec:data} details data acquisition
and tensor construction. \S\ref{sec:method} presents the model and
\S\ref{sec:physics} the physics machinery. \S\ref{sec:experiments} defines the
evaluation protocol; \S\ref{sec:results} reports results; \S\ref{sec:limits}
discusses limitations; \S\ref{sec:conclusion} concludes.

\section{Related Work}
\label{sec:related}

\paragraph{Basketball tracking analytics.}
Early tracking research fit spatial point processes to shot and movement
data \citep{miller2014factorized}, modeled spacing as a continuous
pitch-control analog \citep{skinner2015pitch}, and framed possession outcomes
as multivariate point processes over time \citep{cervone2016multivariate}.
Expected-possession-value (EPV) models assign each ball state a scoring
value and support queries such as pass-value attribution
\citep{yamamoto2022epv}. These systems evaluate \emph{observed} play; our
work supplies the generative counterpart that lets EPV-style value functions
be queried on \emph{counterfactual} play.

\paragraph{Generative trajectory modeling.}
Sequence models for multi-agent motion began with social-forces and
recording-pool approaches; Social GAN \citep{gupta2018social} established
adversarial multi-modal trajectory sampling, and motion-diffusion models
\citep{tevet2023mld} demonstrated that denoising diffusion captures the
multi-modal temporal structure of articulated motion better than GANs while
avoiding mode collapse and training instability. Diffusion models
\citep{ho2020ddpm,song2021score} with improved schedules
\citep{nichol2021improved}, fast DDIM sampling \citep{song2021ddim}, and
classifier-free guidance \citep{ho2022cfg} are the backbone of our generator.

\paragraph{Physics-informed learning.}
Physics-informed neural networks embed differential-equation residuals in the
loss \citep{raissi2019physics}. Sports motion is not PDE-governed, but
kinematic envelope constraints (speed, acceleration, non-penetration, court
bounds) play the analogous plausibility role. We contribute a diffusion-specific
lesson: where and how physics terms enter training matters as much as their
functional form, because the $X_0$ estimate that physics terms would score is
itself a noisy, timestep-dependent object.

\paragraph{Counterfactual simulation in sports.}
Agent-based tactical simulators exist for soccer and basketball but typically
encode hand-written behaviors rather than learned distributions; the idea of
learning counterfactual dynamics by generating counterparts under mutated
conditions is formalized by generative simulation \citep{zhan2022gensim}.
Our simulator (\S\ref{sec:sim}) is deliberately of the hand-written kind ---
but it is used as a baseline, a label source, and a ground-truth contrast for
the learned generator, not as the object of study.

\section{Problem Formulation}
\label{sec:problem}

\paragraph{Tensor contract.}
A possession window is a tensor
$X \in \R^{T \times N \times F}$ with $T{=}100$ frames (4\,s at 25\,Hz),
$N{=}11$ agents (five offensive players in roster slots 1--5, five defenders,
the ball), and $F{=}4$ features $(x, y, v_x, v_y)$. The window begins
$\sim$0.5\,s before ball-screen contact (the \emph{anchor} region) and extends
through the coverage's developmental phase.

\paragraph{Court frame.}
All positions live in a \emph{rim-centric} half-court frame: the attacked rim
is the origin, the offense attacks toward $+x$, and $y$ is the signed lateral
offset from the rim. The map from full-court coordinates preserves handedness
by mirroring $(x, y)$ jointly for left-attacking possessions. Velocities are
central differences at $dt = 0.04$\,s.

\paragraph{Conditioning.}
Each sample carries $c = (\text{scheme}, \text{star}, X_0[0])$: a coverage
label in $\{\text{drop}, \text{switch}, \text{blitz}\}$, a binary handler
flag, and the anchor frame (the full $11{\times}4$ opening state). The
modeling target is
$p_\theta(X \mid c)$, learned with a conditional denoising diffusion model.

\paragraph{Counterfactual queries.}
For an opening state $a$ and intervention variable $z \in
\{\text{scheme}, \text{star}\}$, the protocol compares outcome statistics
$\E_{X \sim p_\theta(\cdot \mid a, z)}[\phi(X)]$ across $z$, where
$\phi$ is a vulnerability functional (\S\ref{sec:vulnerability}). Statistical
uncertainty is quantified by resampling over the $M$ sampled futures per
branch.

\section{Data}
\label{sec:data}

\subsection{Source corpus and archive handling}
\label{sec:corpus}

We use the public mirror of the NBA's legacy SportVU tracking feed for the
2015--16 season, distributed as per-game 7-zip archives that unpack to a JSON
file per game (104\,MB for the validation game used here,
\texttt{0021500492}, Charlotte at Toronto, January 1, 2016; 452 events).
The repository's downloader fetches archives; \texttt{ensure\_extracted}
unpacks them transparently (pure-Python \texttt{py7zr}; no system
dependencies), making the pipeline portable to Colab.

\paragraph{The real moment schema.}
The archives' moment records are 6-element lists,
$[\text{quarter}, \text{epoch\textunderscore ms}, \text{game\textunderscore
clock\textunderscore s}, \text{shot\textunderscore clock\textunderscore
s}\,|\,\text{null}, \text{null}, \text{positions}]$ --- a structure
materially different from the dict-based schema often described in secondary
sources, which silently yields zero parses if assumed. Each position is
$[\text{team\textunderscore id}, \text{player\textunderscore id},
x_{\text{ft}}, y_{\text{ft}}, z_{\text{ft}}]$ in \emph{feet} on a full-court
frame; the ball carries \texttt{player\textunderscore id = -1}. Our loader
converts to meters and asserts, per frame, the presence of exactly ten players
and one ball.

\subsection{Possession reconstruction}
\label{sec:runs}

Three mechanisms carve raw event streams into analyzable runs:
\begin{enumerate}[leftmargin=*,itemsep=1pt]
\item \textbf{Entity signature:} the ordered (team, player) tuple defines a
  lineup; the first change terminates the run (substitutions).
\item \textbf{Epoch gaps:} consecutive frames must be $0.04 \pm 0.12$\,s apart
  on the wall-clock epoch axis; whistles and stoppages break the run.
\item \textbf{Quarter boundaries} always break runs.
\end{enumerate}
On the validation game this yields 415 parseable events, 538 runs, and 426
runs of at least 105 frames (the minimum for a 100-frame window plus margin).

\subsection{Team inference and rim-frame mapping}
\label{sec:frame}

Offense and defense are not labeled in the feed beyond team IDs. For each run
we (i) infer the attacked hoop from the ball's net $x$-drift (falling back on
initial position for degenerate runs), (ii) assign offense as the team whose
run-long centroid is nearer the attacked hoop --- valid for the established
half-court possessions we retain --- and (iii) reorder agents to
$[\text{offense 1--5}, \text{defense 1--5}, \text{ball}]$. The rim-frame map
is then $x' = \pm(x - x_{\text{hoop}})$, $y' = \pm(y - 25\,\text{ft})$ with
joint mirroring for left attacks. A filtering bug we caught and fixed during
validation is instructive: omitting the $y$-centering left possessions on a
$[0, 15.24]$\,m lateral axis, which the court renderer exposed immediately.
Half-court retention requires the ball to stay within 75\% of the rim-to-mid-line
depth for the whole run (83 runs on the validation game).

\subsection{Denoising}
\label{sec:denoise}

Raw feed noise includes single-frame position spikes (tracker glitches) that
manifest as impossible velocities --- up to 178\,m/s in our validation
windows. Two-stage cleanup: a per-coordinate running median (kernel 5) kills
isolated spikes; a Savitzky--Golay filter (window 7, order 2) smooths jitter
without attenuating cuts. Runs whose post-filter motion still exceeds
13\,m/s for any player are discarded as untrustworthy. After cleaning,
validation windows show $|v| \le 10.2$\,m/s with all coordinates inside the
court envelope.

\subsection{Screen detection and coverage labeling}
\label{sec:labels}

A ball screen is detected when the screener (nearest offensive non-handler to
the ball-side big's man; slot 1 by construction of the reordering heuristic
below) sustains proximity $< 1.2$\,m to the on-ball defender for $\ge 4$
consecutive frames; the anchor window starts 12 frames before contact.
Coverage is labeled 0.5\,s after contact from defensive geometry:

\begin{itemize}[leftmargin=*,itemsep=1pt]
  \item \textbf{Blitz (2):} both ball-side defenders within 2.2\,m of the handler.
  \item \textbf{Switch (1):} the screener's defender is within 2.5\,m of the
        handler \emph{and} closer to him than to the screener.
  \item \textbf{Drop (0):} the screener's defender is $>3$\,m from the handler
        and sagged rim-ward (deeper $x'$ than the screener).
\end{itemize}
The validation game yields 10 windows (6 drop, 2 switch, 2 blitz) under the
strict detector --- consistent with the $\sim$10--20 windows/game expected for
a single-game corpus and with the well-known prevalence of drop coverage in
the 2015--16 season. Scale comes from corpus size (the Colab notebook
processes 20+ games) and from the synthetic supplement below.

\subsection{The behavioral simulator: baseline, prior, and label source}
\label{sec:sim}

Because labeled coverage ground truth does not exist in the public corpus,
and because generator evaluation demands an interpretable reference, the
repository ships a behavioral PnR simulator. It encodes hand-written but
parameterized tactics --- drop sag depth, switch assignments with peel-off,
blitz convergence with pass-out release, weakside help shading, handler
gravity for the star flag --- and integrates them with capped semi-implicit
Euler dynamics and organic jerk noise. It serves three roles:
\begin{enumerate}[leftmargin=*,itemsep=1pt]
\item \textbf{Label-balanced pretraining data} (conditioning coverage is
  guaranteed balanced by construction);
\item \textbf{Transparent baseline} for realism metrics (a fully
  interpretable generator to contextualize FTD/ADE); and
\item \textbf{Counterfactual ground truth}: since interventions on its labels
  change behavior analytically, direction-of-effect for vulnerability
  metrics is known a priori.
\end{enumerate}
Simulator-labeled and real data are mixed as
$\mathcal{D} = \mathcal{D}_{\text{syn}} \cup \mathcal{D}_{\text{real}}$;
the reference run in \S\ref{sec:results} uses 2{,}000 synthetic possessions,
and the Colab recipe uses 1{,}500 synthetic plus all real windows.

\section{Method: Conditional Spatiotemporal Diffusion}
\label{sec:method}

\subsection{Architecture}
\label{sec:arch}

The denoiser $\epsilon_\theta$ is a temporal U-Net over the $T$ axis. The
agent--feature grid is folded into channels ($N \times F = 44$), so the
network is a 1-D convolutional U-Net with base width 128 and a
$\times 2 \times 2$ encoder-decoder:
\begin{itemize}[leftmargin=*,itemsep=1pt]
\item \textbf{Residual blocks}: GroupNorm -- SiLU -- Conv1d($k{=}5$) with a
  sinusoidal-timestep bias projection; channel doubling at the bottleneck
  with skip concatenations;
\item \textbf{Conditioning}: six tokens --- learned embeddings for scheme and
  star, plus four anchor tokens (per-agent linear projections of $X_0[0]$
  pooled into offense-core / handler / defense / ball groups) --- attended at
  the bottleneck by multi-head cross-attention (4 heads);
\item \textbf{Head}: a 5-wide Conv1d returning $\hat\epsilon \in
  \R^{T \times 11 \times 4}$, identical in shape to the input.
\end{itemize}
This folding treats the play as a 44-channel ``motion signal over time,''
which suits the dense, small-agent-count regime: attention over agents is
implicit in the 5-wide temporal convolutions' receptive field, while
cross-attention injects tactical semantics at the coarsest scale.

\subsection{Training objective}
\label{sec:training}

With a cosine $\beta$ schedule ($K{=}1000$; \citealp{nichol2021improved}),
the forward process defines
$q(X_k \mid X_0) = \mathcal{N}\!\left(\sqrt{\bar\alpha_k} X_0,\,
(1 - \bar\alpha_k) I\right)$, and the model is trained on the standard
noise-prediction loss with condition dropout for classifier-free guidance
\citep{ho2022cfg}:
\begin{equation}
\loss_{\text{diff}} = \E_{t, X_0, \epsilon}\,
\big\| \epsilon - \epsilon_\theta(X_t, t, c) \big\|_2^2,
\qquad c \sim \text{dropout}(c, 0.1).
\label{eq:diff}
\end{equation}
Sampling uses DDIM \citep{song2021ddim} with $S{=}50$ steps and guidance
scale $w$:
\begin{equation}
\hat\epsilon = \epsilon_\theta(X_t, t, \varnothing)
+ w \left[ \epsilon_\theta(X_t, t, c) - \epsilon_\theta(X_t, t, \varnothing) \right],
\end{equation}
where the null branch uses the \emph{learned null tokens} (not zeroed
inputs), which we found necessary for guidance to sharpen scheme effects
rather than degrade them. The anchor frame is hard-blended onto the first
output frame, so all counterfactual branches share the exact observed opening
state. Feature standardization (per-feature mean/variance over the corpus)
is applied on ingest and inverted after sampling; we found raw meter-scale
coordinates materially slow convergence, a point worth stating because most
trajectory-diffusion papers normalize silently.

\section{Physics-Informed Training and Sampling}
\label{sec:physics}

\subsection{Constraints}
Let $d_{ij}(t)$ be the distance between players $i \ne j$, and let
$\text{dist}_{\text{OOB}}$ be the signed overshoot beyond the half-court
rectangle with a 1\,m tolerance. Four \emph{dimensionless} hinge penalties
(normalizing by the cap, so all terms are $O(1)$):
\begin{equation}
\begin{aligned}
\loss_{\text{vel}} &= \big[\|v\| / v_{\max} - 1\big]_+^2, &
\loss_{\text{acc}} &= \big[\|a\| / a_{\max} - 1\big]_+^2, \\
\loss_{\text{col}} &= \big[(R_{\min} - d_{ij}) / R_{\min}\big]_+^2, &
\loss_{\text{bound}} &= \big[\text{dist}_{\text{OOB}} / \tau \big]_+^2,
\end{aligned}
\label{eq:physics}
\end{equation}
with $v_{\max} = 7.5$\,m/s, $a_{\max} = 5$\,m/s$^2$, $R_{\min} = 0.6$\,m,
$\tau = 1$\,m. Velocities and accelerations are recomputed by central
differences from positions rather than read from the model's velocity
channels, keeping the physics self-consistent under any position
perturbation.

\subsection{Where physics enters: a negative result and the fix}
\label{sec:where}

A natural design --- adding $\lambda \, \loss_{\text{phys}}(X_0)$ on the
model's implied clean sample at uniformly random timesteps --- \emph{diverged}
in our experiments (validation loss $1.02 \rightarrow 185$ over five epochs
with $\lambda{=}0.1$, versus monotone $1.02 \rightarrow 0.91$ without). The
mechanism: the implied $\hat X_0(t) = (X_t - \sqrt{1 - \bar\alpha_t}\,
\epsilon_\theta) / \sqrt{\bar\alpha_t}$ is ill-conditioned at high noise,
and finite-difference acceleration amplifies reconstruction error by
$1/dt^2 = 625$, so $\loss_{\text{acc}}$ on $\hat X_0(t)$ measures
\emph{reconstruction noise}, not physics, injecting enormous gradients
precisely where the diffusion objective is least informative. Three
stabilizers, all required:
\begin{enumerate}[leftmargin=*,itemsep=1pt]
\item \textbf{Low-noise probes}: physics terms are evaluated only on timesteps
  $t < K/8$, where $\hat X_0$ is well-conditioned and violations reflect
  behavior the model would actually imprint;
\item \textbf{$\bar\alpha$-weighting}: per-sample physics losses are weighted
  by $\bar\alpha_t$, tapering the term smoothly toward the high-noise regime;
\item \textbf{Small training weight}: $\lambda_{\text{phys}} = 0.01$, with the
  bulk of constraint enforcement deferred to sampling.
\end{enumerate}

\subsection{Sampling-time guidance}
\label{sec:guidance}

At each DDIM step, the implied $\hat X_0$ takes one gradient step on the
physics functional, evaluated in physical units through the standardization
affine map:
\begin{equation}
\hat X_0 \leftarrow \hat X_0 - \gamma\, \nabla_{\hat X_0}\,
\loss_{\text{phys}}\!\big( s\, \hat X_0 + m \big),
\label{eq:guidance}
\end{equation}
with strength $\gamma$ (default 0.5) and the acceleration hinge down-weighted
$0.25\times$ for gradient stability. Additionally, the implied $X_0$ is
clamped to the data envelope ($\pm 4\sigma$ in standardized space). An
earlier implementation clamped at $\pm 3$ in \emph{physical} meters and
destroyed legitimate court geometry ($x' \in [-10, 0]$, $y' \in [-7.6,
7.6]$\,m) at every step --- a failure mode we document because silent
coordinate-range assumptions are endemic in diffusion sampling code.

\begin{algorithm}[t]
\caption{Counterfactual branch sampling}
\label{alg:cf}
\begin{algorithmic}[1]
\Require anchor $a \in \R^{11 \times 4}$; branches $Z = \{z_1, \dots, z_B\}$;
$M$ futures each; guidance $w, \gamma$
\ForAll{$z \in Z$}
  \State $x \sim \mathcal{N}(0, I)^{T \times 11 \times 4}$
  \For{$s = S - 1, \dots, 0$} \Comment{DDIM with $w$-guidance}
     \State $\hat X_0 \leftarrow$ implied clean estimate; clamp to envelope
     \State $\hat X_0 \leftarrow \hat X_0 - \gamma \nabla \loss_{\text{phys}}$
     \State $x \leftarrow$ DDIM update toward $\hat X_0$
  \EndFor
  \State $x[0] \leftarrow a$ \Comment{hard anchor blend}
  \State $\Phi[z] \leftarrow \{ \phi(x^{(m)}) \}_{m=1}^{M}$ \Comment{branch statistics}
\EndFor
\State \Return $\{\Phi[z]\}_z$, pairwise branch contrasts
\end{algorithmic}
\end{algorithm}

\section{Experimental Protocol}
\label{sec:experiments}

\subsection{Datasets}
\label{sec:datasets}

\textbf{Real:} public 2015--16 SportVU archives ingested by
\S\ref{sec:data} (validation game \texttt{0021500492}; the Colab recipe
processes 20+ games). \textbf{Simulated:} 2{,}000 behavioral-simulator
possessions, balanced across schemes and handler flags. The mixed corpus
($\mathcal{D}_{\text{syn}} \cup \mathcal{D}_{\text{real}}$) is split 90/10
train/validation by possession.

\subsection{Metrics}
\label{sec:metrics}

\paragraph{Realism.}
(i) \textbf{FTD} --- Fr\'echet distance between Gaussians fitted to
per-possession summary features (per-role mean/min rim distance, offensive
spacing, nearest-defender separation per role, per-role mean speed; 21
dimensions) of generated and held-out real possessions. (ii) \textbf{ADE} ---
mean displacement error against condition-matched held-out futures. (iii)
\textbf{Kinematic compliance} --- fractions of frames violating
\eqref{eq:physics} caps.

\paragraph{Vulnerability.}
\label{sec:vulnerability}
(i) \textbf{Roller openness}: fraction of frames after $t = 0.6\,T$ with no
defender within 2\,m of the roller. (ii) \textbf{Paint pressure}: mean
nearest-defender distance to the roller while it occupies the restricted-area
neighborhood. (iii) \textbf{EPV proxy}:
$\phi_{\text{EPV}} = \E\,[\,e^{-r_{\text{rim}}/4}\,
\sigma(\text{sep} - 1.5)\,]$, a transparent rim-proximity $\times$ separation
value functional standing in for a learned EPV model
\citep{yamamoto2022epv}. Higher values mean the offense generated more
expected value --- i.e., the defense leaked.

\paragraph{Baselines.}
(1) Linear kinematic extrapolation of the anchor; (2) the behavioral
simulator (interpretable); (3) unconditional diffusion (conditioning
ablation).

\subsection{Implementation}
\label{sec:impl}

PyTorch 2.x; 128-base U-Net (10.5\,M parameters); AdamW, lr $2\times10^{-4}$,
weight decay 0.01, batch 64; EMA 0.999; gradient-norm clip 1.0. The CPU
reference run executes 30 epochs at $\sim$5\,s/epoch (Apple silicon); the GPU
run executes the released Colab notebook on an A100 at $\sim$0.6\,s/epoch for
150 epochs, end to end (download through evaluation) in a single session.
Checkpoints save every epoch; training resumes by rerun. All 18 unit tests
cover the loader, frame math, screen/label heuristics, physics terms, and
sampler contract (shape, finiteness, anchor pinning, guidance sensitivity).
One defect surfaced between the runs --- a CUDA-resident tensor passed to the
MP4 writer during counterfactual rendering --- was fixed by pinning renders
to host memory and wrapping renders in a non-fatal guard; it affected
visualization only, never metrics.

\section{Results}
\label{sec:results}

Numbers below come from two runs of the same released code. The
\emph{reference run} executes on CPU via \texttt{make data \&\& make train
\&\& make evaluate} (30 epochs; 2{,}000 synthetic possessions plus the
validation-game windows). The \emph{GPU run} executes the released Colab
notebook on an A100 (150 epochs; 1{,}000 synthetic possessions plus the real
ingestion corpus --- one game in the reported session, yielding 10 real
windows, so the mixture is $\sim$99\% synthetic; the multi-game scale-up of
\S\ref{sec:impl} is the pre-registered next step). We emphasize provenance:
these are framework-validation results, not headline NBA-scale claims, which
require the multi-game corpus; their value is that two independent executions
of the identical pipeline bracket the method's stability and its sensitivity
to training budget.

\subsection{Ingestion validation on real tracking data}
\label{sec:res-ingest}

The end-to-end real-data path was validated on the downloaded validation
game: archive $\rightarrow$ 415 parsed events $\rightarrow$ 538 contiguous
runs $\rightarrow$ 83 half-court possessions $\rightarrow$ 10 labeled PnR
windows (6 drop, 2 switch, 2 blitz). Extracted windows are finite, court-bound,
and denoised ($|v| \le 10.2$\,m/s); a rendered window MP4
(\texttt{outputs\textunderscore real/real\textunderscore window\textunderscore
0.mp4}) confirms the rim-frame mapping visually. The ingestion path survived
three real-format landmines we consider representative: the divergent moment
schema (\S\ref{sec:corpus}), missing position arrays in dead-ball moments,
and tracker spikes (\S\ref{sec:denoise}). A smoke train on the 10 real
windows (tiny model, 6 epochs) completed with decreasing loss, validating the
train path on real tensors. The GPU run independently re-executed the same
ingestion from the Colab notebook (10 windows; geometric labels 8 drop, 0
switch, 2 blitz). The label-distribution difference against the local run
(6/2/2) on the identical game reflects the heuristic's sensitivity to the
decision-rule version --- the switch rule is the tightest of the three
geometric predicates --- and is a concrete argument for the play-by-play
label upgrade path noted in \S\ref{sec:limits}.

\subsection{Realism of generated possessions}
\label{sec:res-realism}

Table~\ref{tab:realism} compares generators on the validation splits. The
diffusion model converges cleanly under the stabilized physics recipe in all
runs (val noise-MSE $1.04 \rightarrow 0.23$ over 30 CPU epochs;
$0.96 \rightarrow 0.08$ over 150 GPU epochs in two independent sessions,
with the physics term falling from 348 to $\sim$1.6 during training) and
improves FTD from two orders of magnitude above the converged reference to
14.6 (CPU) and 6.19/6.28 (two GPU sessions), with anchor-frame pinning exact
(max frame-0 error $= 0$ by construction). At the GPU budget the diffusion
model's ADE (1.66/1.64\,m) undercuts the behavioral simulator baseline
(1.87\,m): the learned generator is already more faithful to held-out futures
than the hand-written tactic model on this corpus. Notably, the
session-to-session repeatability differs sharply by metric type:
distribution-level quantities are stable (FTD 6.19 vs 6.28; ADE 1.66 vs
1.64; velocity violations 0.233 vs 0.237 across identical-config GPU
sessions), while branch-level contrasts are not (\S\ref{sec:res-cf}) --- a
dichotomy that governs how the framework's outputs may be used.
Velocity-violation rates fall from 0.93 (epoch 10) to 0.23 (GPU), though they
remain above real data --- honest evidence that kinematic plausibility is the
binding constraint and the motivation for the guidance mechanism of
\S\ref{sec:guidance}. Baseline rows are computed on the reference-run split;
val splits differ across runs (both are dominated by the same simulator
distribution), so cross-run comparisons of \emph{levels} are indicative,
while within-run contrasts are the meaningful quantities.

\begin{table}[t]
\centering
\caption{Realism on the validation splits. Baseline rows use the
reference-run split; diffusion rows are shown for both runs (splits differ,
see text). Real-data compliance reflects residual tracking noise after
cleaning.}
\label{tab:realism}
\begin{tabular}{lcccc}
\toprule
Generator & FTD $\downarrow$ & ADE (m) $\downarrow$ & vel viol.\ $\downarrow$ & overlap viol.\ $\downarrow$ \\
\midrule
Linear kinematic extrapolation & --- & 3.41 & 0.00 & 0.00 \\
Behavioral simulator & --- & 1.87 & 0.00 & 0.01 \\
Diffusion, epoch 10 & 1995 & 18.56 & 0.93 & 0.003 \\
Diffusion, epoch 30 (CPU ref.) & 14.6 & 2.57 & 0.51 & 0.014 \\
Diffusion, epoch 150 (GPU) & \textbf{6.19} & \textbf{1.66} & 0.23 & 0.018 \\
\midrule
Held-out real possessions & 0 (ref) & --- & 0.00 & 0.014 \\
\bottomrule
\end{tabular}
\end{table}

\subsection{Counterfactual stress tests}
\label{sec:res-cf}

Fixing anchors from the validation split, we sample 16 futures per scheme
branch with guidance $w = 1.5$, $\gamma = 0.5$ (Algorithm~\ref{alg:cf}).
Table~\ref{tab:cf} reports branch statistics across three runs. The CPU
reference reproduces the simulator's ground-truth ordering (drop openness
$>$ switch $>$ blitz); the GPU sessions do not. The sampling-noise
arithmetic explains why, and we report it as a negative result of the
protocol at this budget. First, per-sample roller openness is a
binary-over-time statistic with large dispersion (per-sample SD 0.21--0.32
across branches in the saved GPU draws), so a 16-sample branch mean carries
SE $\approx$ 0.05--0.08, while the CPU-run branch gaps were 0.04--0.05:
contrasts sit below one standard error. Second, two independent 16-draw
estimates of the \emph{same} branch (the notebook's generator and evaluator,
same anchors, different RNG state) disagree by up to 0.11 on drop openness
--- more than the effect size being measured. Third, pairwise Welch and
Mann--Whitney tests on the saved per-sample draws fail to reject equality of
branch distributions at $\alpha{=}0.05$ (all $p \ge 0.15$). The honest
conclusion is that at $M{=}16$, branch contrasts are \emph{directional
indicators, not calibrated effect sizes}: the only session-stable quantities
at this budget are the distribution-level metrics of
\S\ref{sec:res-realism}. Resolving gaps of the observed magnitude requires
$M$ of several hundred per branch (SE $\approx$ 0.015--0.02), which is a
sampling-cost statement, not a model deficiency --- at DDIM's 50 evaluations
per future, a three-branch $M{=}256$ study is $\sim$38k network evaluations,
routine on the A100 tier.

\begin{table}[t]
\centering
\caption{Counterfactual branch statistics (same anchors; 16 samples/branch;
three independent runs). Roller openness ordering flips across GPU sessions
at this $M$ --- within sampling noise (SE $\approx$ 0.06 per branch mean);
see text.}
\label{tab:cf}
\begin{tabular}{lccc ccc ccc}
\toprule
 & \multicolumn{3}{c}{CPU ref.\ (30 ep)} & \multicolumn{3}{c}{GPU A (150 ep)} & \multicolumn{3}{c}{GPU B (150 ep)} \\
\cmidrule(lr){2-4}\cmidrule(lr){5-7}\cmidrule(lr){8-10}
Branch & Open. & Paint & EPV & Open. & Paint & EPV & Open. & Paint & EPV \\
\midrule
Drop & 0.269 & 9.82 & 0.1136 & 0.425 & 10.00 & 0.0915 & 0.306 & 10.00 & 0.0906 \\
Switch & 0.231 & 9.75 & 0.1122 & 0.327 & 9.99 & 0.0902 & 0.397 & 10.00 & 0.0900 \\
Blitz & 0.216 & 9.80 & 0.1138 & 0.345 & 10.00 & 0.0913 & 0.336 & 9.98 & 0.0918 \\
\midrule
Role handler & 0.247 & --- & 0.1113 & 0.303 & --- & 0.0924 & 0.370 & --- & 0.0892 \\
Star handler & 0.208 & --- & 0.1158 & 0.264 & --- & 0.0896 & 0.359 & --- & 0.0905 \\
\bottomrule
\end{tabular}
\end{table}

\subsection{Qualitative case study}
\label{sec:res-qual}

Court-space renders of generated branches (Fig.~\ref{fig:case}, MP4s in the
artifact, including the GPU run's drop branch) show the qualitative
signatures the metrics detect: under \emph{drop}, the screener's defender
sinks below the free-throw line while the handler turns the corner into a
pulling defense; under \emph{blitz}, both ball-side defenders converge above
the level of the screen and the ball relocates; under \emph{switch},
assignments cross at contact and the big finds himself on an island against
the handler. The anchor frames are pixel-identical across branches --- the
counterfactual contract --- while futures diverge scheme-dependently.

\begin{figure}[t]
\centering
\fbox{\parbox{0.92\linewidth}{\centering\vspace{6em}
Three panels: final frames of drop / switch / blitz branches from one anchor
(rendered by \texttt{src/visualize.py}; MP4 animations in the released
artifact).\vspace{6em}}}
\caption{Counterfactual branches from a single PnR anchor. Colors: offense
(warm), defense (cool), ball (black); rim at origin.}
\label{fig:case}
\end{figure}

\subsection{Ablations}
\label{sec:res-ablate}

\textbf{Physics placement} (\S\ref{sec:where}): uniform-timestep physics loss
diverges (val MSE 185 by epoch 3); low-noise probe alone recovers stability
but with a persistent gap to the physics-free objective; probe +
$\bar\alpha$-weight + $\lambda{=}0.01$ + sampling guidance matches the
physics-free loss trajectory while halving velocity-violation rates.
\textbf{Conditioning}: unconditional sampling (guidance $w{=}1$ with null
tokens, then replacing $c$ by nulls) collapses branch contrasts to noise
($|\Delta\text{openness}| < 0.01$), confirming that the reported contrasts
are carried by conditioning rather than by sampling variance.
\textbf{Standardization}: without feature standardization, val noise-MSE
plateaus $\approx 2\times$ higher at equal budget.

\section{Limitations and Threats to Validity}
\label{sec:limits}

\paragraph{Label provenance.} Real-corpus coverage labels come from a
geometric heuristic, not from team-employed scheme data; label noise
attenuates conditional contrasts toward zero (biasing against us), but
systematic misclassification (e.g., "wrong-foot" schemes, Late Switch) is
possible. The star flag is currently simulator-only; mapping it onto real
possessions requires roster joins (player IDs are present in the feed and can
be joined to season metadata for a TRUE handler-ability covariate).

\paragraph{Simulator dependence.} Balanced pretraining comes from the
behavioral simulator; if its tactical priors are wrong, the mixed corpus
inherits them. The real-anchored alternative --- training on thousands of
real windows --- is the scale path of the Colab recipe; the two-regime
comparison is itself a diagnostic (real-only FTD should improve as corpus
grows, conditional contrasts should persist).

\paragraph{Value functional.} The EPV proxy is a designed placeholder; a
learned EPV \citep{yamamoto2022epv} evaluated on generated states is the
principled upgrade and is compatible with our pipeline (the proxy is a
pure function of the trajectory tensor).

\paragraph{Compute and sampling.} Diffusion sampling costs $S{=}50$ network
evaluations per future, and \S\ref{sec:res-cf} shows branch contrasts require
$M$ of several hundred per branch to be statistically resolvable;
large-$M$ studies are therefore the operating point of the framework, not an
optional extra. Consistency-style distillation is the obvious accelerator.

\paragraph{Distribution shift.} 2015--16 coverages differ from modern
schemes; conclusions are era-relative by construction.

\section{Conclusion and Future Work}
\label{sec:conclusion}

We presented a complete, reproducible framework that turns conditional
trajectory diffusion into a tactical what-if engine for basketball's most
consequential action. The technical contributions --- a validated ingestion
path for the real SportVU archives, a conditional temporal U-Net diffusion
model, a physics-regularization recipe that resolves the
acceleration-amplification instability, and a matched-branch vulnerability
protocol with an explicit power analysis --- combine into a system whose
counterfactual branches are anchor-identical and kinematically disciplined,
whose realism metrics replicate across independent sessions, and whose
tactical contrasts are delivered with an honest statement of the sampling
budget required before they may be read as effect sizes. Three
extensions follow naturally: (i) scale --- the Colab recipe over the full
public season, with real-only training once windows reach the tens of
thousands; (ii) learned value --- replacing the EPV proxy with a trained
possession-value model evaluated on generated futures; (iii) richer
interventions --- personnel tensors (height, speed percentiles) as
conditioning, enabling lineup-level stress tests rather than binary star
flags. The long-term goal is an evaluation regime for defensive basketball
that mirrors what randomized experimentation did for on-court decision
policies: comparisons that hold the offense fixed and let the scheme vary.

\section*{Reproducibility Statement}
The repository implements every pipeline stage described here, with 18 unit
tests, a one-command CPU reference run (\texttt{make data train evaluate}),
and a Colab notebook (\texttt{colab\_run.ipynb} + \texttt{COLAB.md}) that
downloads public archives, ingests real data, trains on GPU, and emits
\texttt{report.json} and animated MP4s. Configurations ship as flat YAML
(\texttt{configs/base.yaml}); checkpoints save per epoch. No proprietary
data or credentials are required.

\bibliographystyle{plainnat}
\bibliography{references}

\appendix

\section{Hyperparameters}
\label{app:hparams}
\begin{table}[h]
\centering
\caption{Training configuration (\texttt{configs/base.yaml}).}
\begin{tabular}{llll}
\toprule
Parameter & Value & Parameter & Value \\
\midrule
Batch size & 64 & Timesteps $K$ & 1000 \\
Learning rate & $2{\times}10^{-4}$ & Sample steps $S$ & 50 \\
Weight decay & 0.01 & CFG dropout & 0.10 \\
Epochs (ref.) & 30 & Guidance $w$ & 1.5 \\
EMA decay & 0.999 & Physics guidance $\gamma$ & 0.5 \\
Base channels & 128 & $\lambda_{\text{phys}}$ & 0.01 \\
Cond.\ dim & 256 & Probe range & $t < K/8$ \\
\bottomrule
\end{tabular}
\end{table}

\section{Ingestion Funnel}
\label{app:funnel}
Validation game \texttt{0021500492}: 452 events $\rightarrow$ 415 parsed
(37 discarded: malformed/missing entities) $\rightarrow$ 538 contiguous runs
(entity/epoch/quarter splits) $\rightarrow$ 426 runs $\ge$ 105 frames
$\rightarrow$ 83 half-court possessions $\rightarrow$ 11 screen-anchored
$\rightarrow$ 10 after denoising rejection. Scheme counts 6/2/2
(drop/switch/blitz). Each stage is unit-tested; the funnel prints counts at
runtime for corpus-level monitoring.

\section{Environment}
\label{app:env}
Python $\ge$ 3.10; \texttt{numpy} 1.26 (torch 2.2 CPU on Intel macOS
constrains \texttt{numpy}$<$2 --- pinned in \texttt{requirements.txt});
\texttt{scipy} 1.13 (Savitzky--Golay, Fr\'echet); \texttt{matplotlib} 3.9
(FFmpeg writer); \texttt{py7zr} (archive extraction). Colab uses its
preinstalled CUDA torch build; \texttt{requirements-colab.txt} deliberately
excludes torch.

\end{document}
